% ---------------------------------------------------------------------------
% preface.tex
% ---------------------------------------------------------------------------
\chapter*{Preface}
\addcontentsline{toc}{chapter}{Preface}
\markboth{Preface}{Preface}

Sooner or later you will need to run a simulation without a QuestaSim
licence in front of you, or with a golden model already written in Python
that you would rather not reimplement in SystemVerilog. This booklet is
about both situations.

This is the \textbf{Simulation Cookbook}, one of four companions covering
\tool{SystemVerilog for VHDL Designers} in full. It assumes the design and
verification material of \refdesignbook and \refverificationbook, and adds
the tools that make a testbench runnable without a commercial licence, plus
the four ways of driving one from Python.

\section*{What this booklet is}

\Cref{ch:verilator} is about \tool{Verilator}, which is not an event
simulator and should not be treated as one: it compiles a language subset to
C++, it is two-state, and it has no route to UVM, but it is also the fastest
thing you will run on a laptop and the best free linter for SystemVerilog by
a wide margin. The chapter compares it directly with \tool{QuestaSim},
feature by feature, so you know before you start which of the two a given
testbench needs. \Cref{ch:python} covers the four ways of connecting Python
to a testbench, from a bare socket to \tool{cocotb}, which lets the HDL
testbench disappear entirely.

\section*{What this booklet is not}

It is not a verification methodology text; the class-based testbench and UVM
that these tools ultimately run are covered in \refverificationbook, and the
design-side language they compile is covered in \refdesignbook. Where this
booklet needs a design to simulate, it uses the CORDIC rotator developed in
full in \refexamplebook.

\section*{The tools}

\Cref{app:tooling} lists exactly what to install, on Linux, macOS or WSL2,
and \cref{app:simcommands} gathers every command line used in this booklet
onto one page. Nothing in \cref{ch:verilator} or \cref{ch:python} needs a
licence.
